\documentclass{article}

\textwidth=6.5in
\textheight=8.5in
\voffset=-54pt
\oddsidemargin=0in
\evensidemargin=0in
\setlength{\parskip}{8pt}
\setlength{\parindent}{0pt}

\usepackage[normalem]{ulem}

\usepackage{amsmath, amssymb}
\usepackage{ifthen}

\usepackage[inline]{enumitem}
\setlist{topsep=0pt}

\usepackage{xcolor}
\usepackage[pdftex]{graphicx}

\usepackage{siunitx}
\sisetup{per-mode=fraction}

% change hyperref options
\PassOptionsToPackage{hyphens}{url}
\usepackage[bookmarks,colorlinks,linkcolor=black,citecolor=blue,pdfstartview=FitH,urlcolor=black]{hyperref}
\hypersetup{pdfpagemode=UseNone}

\usepackage{graphicx}
\usepackage{multicol}
\usepackage{framed}
\usepackage{array}
\usepackage{icomma}
\usepackage{diffcoeff}

\usepackage{soul}

\usepackage{pgfplots}
\pgfplotscreateplotcyclelist{mycolor}{%
  black, blue, red, green!70!black,magenta,
  purple, cyan, orange
}
\pgfplotsset{
  cycle list name=mycolor,
  axis lines=middle,
  every axis plot/.style={no marks, thick},
  every axis/.style={
    enlarge x limits=false,
    enlarge y limits=auto,
    xlabel style={at={(current axis.right of origin)},anchor=west},
    ylabel style={at={(current axis.above origin)},anchor=south},
    % extra x and y ticks for asymptotes
    extra x tick style={grid=major, grid style={dashed,black}},
    extra y tick style={grid=major, grid style={dashed,black}},
    /pgf/number format/1000 sep={},
  },
  tick label style = {font=\footnotesize, fill=white, inner sep=0pt},
  xticklabel shift=1pt,    
  scaled ticks=false,
  scaled y ticks=false,
  unbounded coords=jump,
  samples=101,
  minor tick length=0,
  scatter/classes={
    open={mark=*,draw=black,fill=white, mark size=2, thin},
    closed={mark=*,draw=black,fill=black, mark size=2, thin}
  },
  compat=1.18
}

\usepackage{tikz}
\usetikzlibrary{
  % enables the snake paths
  decorations.pathmorphing,%
  % enables bigger arrows
  decorations.markings,%
  decorations.pathreplacing,%
  decorations.text,%
  calc,%
  patterns,%
  shapes.geometric,%
  arrows,
  decorations.shapes,
  positioning,
  plotmarks,
}
\tikzset{>=latex} % sets default arrow type in tikz
\tikzset{font=\small}

\interfootnotelinepenalty=10000
\renewcommand{\thefootnote}{(\arabic{footnote})}

\usepackage{multirow, bigdelim}

\newcommand\iscurrentenumi[1]{%
  \ifthenelse{\equal{#1}{\theenumi}}{}{#1}}
\setenumerate[1]{ref=\#\arabic*}
\setenumerate[2]{ref=\protect\iscurrentenumi{\theenumi}(\alph*)}
\newcommand\enumiiprefix[2]{%
  \ifthenelse{{\equal{#1}{\theenumi}}\AND{\equal{#2}{(\alph{enumii})}}}{}{}%
  \ifthenelse{{\equal{#1}{\theenumi}}\AND{\NOT{\equal{#2}{(\alph{enumii})}}}}{#2}{}%
  \ifthenelse{\equal{#1}{\theenumi}}{}{#1#2}%
}
\setenumerate[3]{ref=\protect\enumiiprefix{\theenumi}{(\alph{enumii})}\roman*}

\definecolor{ideacolor}{rgb}{0.7,0,0}
\newcommand{\ideas}[1]{\color{ideacolor} #1 \color{black}}
\newcommand{\ideacolor}{maroon}

\newcommand{\outside}[1]{
    \fcolorbox{red}{white}{#1}
}

\newcommand{\inside}[1]{
    \fcolorbox{blue}{white}{#1}
}

\newlength{\tipmargin}
\makeatletter
\newcommand{\version}[2]{\ifthenelse{\equal{#1}{Mb}}{#2}{}}

\renewenvironment{framed}{
  \setlength{\tipmargin}{\textwidth-\linewidth}
  \advance\@totalleftmargin-\tipmargin
  \def\FrameCommand{\hspace{\tipmargin}\fbox}%
  \advance\linewidth\tipmargin
  \MakeFramed {\advance\hsize-\width \FrameRestore}%
}
{\endMakeFramed}

\makeatother

\newlength{\highlightwidth}
\newcommand{\highlight}[1]{\setlength{\highlightwidth}{\linewidth}%
  \addtolength{\highlightwidth}{-55pt}%
  \hspace{24pt}\begin{minipage}[t]{\highlightwidth} \setlength{\parskip}{8pt} #1 \end{minipage}\hspace{24pt}}

\definecolor{purple}{rgb}{0.5,0,1.0}

\newcommand{\textbook}[0]{\href{http://isites.harvard.edu/fs/docs/icb.topic1511355.files/fullbook.pdf}{textbook}}

% change to \usepackage[sols]{handout} to see solutions
% change to \usepackage[notes]{handout} to see notes
\usepackage[]{handout}
\renewcommand{\workshopnote}{CA Notes.}    

\begin{document}

\htitle{Worksheet 1 - The Chain Rule}
\begin{problemonly}
    \begin{framed}

You should be able to:

\begin{itemize}[topsep=0pt,partopsep=0pt,parsep=8pt,itemsep=0pt]

\item You should be able to use the Chain Rule to find the derivative of
  the composition of two functions.

\item You should understand why the Chain Rule makes sense in a real-world
  problem (such as {{\#1}} on {the worksheet ``The Chain Rule''}).

\item You should be comfortable looking at a function and determining its
  structure: is it a product?  A composition?  (If so, what is the outer
  function?)

\end{itemize} 

    \end{framed}
\end{problemonly}

\begin{enumerate}
\item 
  \note{I'll use this to try to convince them that the Chain Rule is
    reasonable.}

  A rainwater collection tank is shaped so that 
  when the tank is filled with $v$ liters of rainwater,
  then the height of water in the tank will be $v^2$ mm.
  (The tank narrows at the top).
  \begin{enumerate}
  \item
    \note{Students' first instinct here will probably be to take the
      derivative of $f(v) = v^2$ and plug in $v = 10$.  (Make sure they
      are absolutely solid with the connection between ``instantaneous
      rate of change'' and ``derivative''.)  So, we'll talk about how 
      $f'(10)$ is $\frac{dh}{dv}$ and what this represents.  From an
      intuitive perspective, they should see that the answer to the
      question must depend on the precipitation rate: if it's not
      raining at all, the rate in question is definitely 0. 

      Something that will probably seem pretty new to them is thinking of
      a quantity (height, in this case) as a function of different
      independent variables (volume vs. time).
    }

    \begin{problem}[\vfill]
      During a storm, the tank begins filling up.  At 1 pm, the tank contains
      exactly 10 L of water.  Can you tell what the
      instantaneous rate of change of the height of the water with respect to time
      is at 1 pm, or do you need more information?
    \end{problem}

    \begin{solution}
      We need more information; we aren't told how quickly water is
      entering the tank, so there's no way we could tell the instantaneous
      rate of change of the height of the water with respect to time.  (For
      instance, if the rain had temporarily paused at 1 pm, then
      the instantaneous rate of change of height of water with respect to
      time at that moment would be 0.)
    \end{solution}

  \item

    \note{I'll use this to get the Chain Rule in Leibniz notation:
      $\displaystyle \frac{dh}{dt} = \frac{dh}{dv} \cdot \frac{dv}{dt}$.
      (I won't try to prove anything.) Some students will immediately
      notice the similarity to the dimensional analysis that they learned
      in high school. It's worth bringing this up if none of the students
      does.

      If you can, it's worthwhile to already start using the framework
      we'll use in related rates: identify the rate you want
      ($\frac{dh}{dt}$) and the rate you're told ($\frac{dv}{dt} = 2$). 
    }

    \begin{problem}[\vfill]
      Suppose that you are also told that, at 1 pm, rainwater is
      pouring into the tank at an instantaneous rate of 2 L/s.  Now, can
      you tell what the instantaneous rate of change of the height of the water
      with respect to time is, or do you need more information?
    \end{problem}

    \begin{solution}
      Intuitively, it should feel like we have enough information because
      we're now told how quickly water is entering the tank. Since we know
      that water is entering the tank at a rate of $2$ L/s, to determine the
      how fast the height of the water is increasing, we just need to determine
      the rate at which the height increases with respect to volume.

      We are told that $h=v^2$, so the rate of change of the height of the water 
      in the tank with
      respect to the volume is $\frac{dh}{dv} = 2v$. This means that, when
      there are $10$ L of water in the tank, the $\frac{dh}{dv} = 2 \cdot 10 =
      20$ mm/L.

      Combining the two pieces of information, we can determine that the
      instantaneous rate of change of the height of the water in the tank
      with respect to time is 
      \[\qty{2}{\liter\per\second} \cdot \qty{20}{\milli\meter\per\liter} = 
      \qty{40}{\milli\meter\per\second}.\]

      This demonstrates the following relationship between $\diff{h}{t}$,
      $\diff{h}{v}$, and $\diff{v}{t}$:
      \[
        \frac{dh}{dt} = \frac{dh}{dv} \frac{dv}{dt}
      \]
    \end{solution}

  \end{enumerate}

\item 
  \note{The goal here is to connect {{\#1}} with the Chain
    Rule they saw on the homework. It's important that students are
    comfortable with both Leibniz and Lagrange notation.}

  \begin{problem}[\vfill]
    On your homework, you looked at the derivative of a composition of
    functions, $h(t) = f(g(t))$.  What does this have to do with what we
    figured out in {{\#1}}?
  \end{problem}

  \begin{solution}
    Let's say $v = g(t)$ and $h = f(v)$.  What we learned from
    {{\#1}} was that $\displaystyle \frac{dh}{dt} =
    \frac{dh}{dv} \cdot \frac{dv}{dt}$.  Let's see what this says in terms
    of $f$ and $g$.  First of all, we know that
    \begin{align*}
      h &= f(v) \\
      &= f(g(t)) \text{ since $v = g(t)$} \\
      &= h(t)
    \end{align*}
    So, another way of writing $\frac{dh}{dt}$ is $h'(t)$.    On the other
    hand, $\frac{dh}{dv} = f'(v)$ 
    and $\frac{dv}{dt} = g'(t)$.  So, we can rewrite the equation
    $\displaystyle \frac{dh}{dt} = \frac{dh}{dv} \cdot \frac{dv}{dt}$ as
    \begin{align*}
      h'(t) &= f'(v) g'(t) \\
      \intertext{Since $v = g(t)$,
        this is the same as saying}
      h'(t) &= f'(g(t)) g'(t),
    \end{align*}
    and this is exactly the Chain Rule you looked at in your homework. 
  \end{solution}

\pspace{1in}

\begin{framed}
    \textbf{The Chain Rule:} Given a composition of functions $f(g(t))$, the derivative is:
    $$\dfrac{d}{dt}[f(g(t))]= \solonly{f'(g(t))\cdot g'(t) }
    \hspace{1.5 in} \text{or} \hspace{.5 in} \dfrac{df}{dt} = 
    \solonly{\dfrac{df}{dg} \cdot \dfrac{dg}{dt}}
    $$ 
\end{framed}

\pspace{\newpage}

\item 
  \note{This will hopefully help them see better how we apply the Chain
    Rule in practice.  For example, when we want to differentiate ln of
    something, students sometimes need help seeing that they can use the
    Chain Rule by viewing it as $\ln [ g(x)]$.  You might encourage the
    students to fill this out for themselves first, before checking with a
    neighbor; if you do this, you might also mention that self-quizzing is
    a very effective studying tactic.}

  Fill in the following tables.
  \begin{center}
    \begin{tabular}{| l | l | c | l | l |}
      \multicolumn{2}{c}{Table 1: Derivatives we knew before} &
      \multicolumn{1}{c}{} 
      & \multicolumn{2}{c}{Table 2: Derivatives we know now} \\ 
      \cline{1-2} \cline{4-5} & & & & \\
      $f(x) = x^n$ & $f'(x) = \solonly{nx^{n-1}}$ \hspace{0.4in} &
      \hspace{0.5in} &
      $h(x) = [g(x)]^n$ & $h'(x) = \solonly{n[g(x)]^{n-1}g'(x)}$
      \hspace{0.9in} \\ 
      & & & & \\
      $f(x) = e^{kx}$ & $f'(x) = \solonly{ke^{kx}}$ & & $h(x) = e^{g(x)}$
      & $h'(x) = \solonly{e^{g(x)}g'(x)}$ \\ 
      & & & & \\
      $f(x) = b^{x}$ & $f'(x) = \solonly{\ln(b)b^x}$ & & $h(x) = b^{g(x)}$
      & $h'(x) = \solonly{\ln(b)b^{g(x)}g'(x)}$ \\ 
      & & & & \\
      $f(x) = \ln x$ & $f'(x) = \solonly{\frac{1}{x}}$ & & $h(x) = \ln
      [g(x)]$ & 
      $h'(x) = \solonly{\frac{g'(x)}{g(x)}}$ \\ 
      & & & & \\
      \cline{1-2} \cline{4-5} 
    \end{tabular}
  \end{center}
  
  \note{In class, I write another table:
  \renewcommand{\arraystretch}{1.5}
  \begin{tabular}{| l | l |}
  \hline
  $ y = (\text{mess})^n $ & $ y' = n(\text{mess})^{n-1} (\text{mess})' $ \\
  $ y = e^\text{mess} $ & $ y' = e^\text{mess} (\text{mess})' $\\
  $ y = \ln(\text{mess}) $ & $ y' = \dfrac{1}{\text{mess}}(\text{mess})' $ \\
  \hline
  \end{tabular} and then I comment that I always try to use log rules
  to simplify the mess because making extra work for oneself is NOT a virtue! For example, $\ln(x^3/(2x+7))$ could be done using the chain rule and quotient rule, but it's so much simpler to write this as $3\ln x - \ln(2x+7)$ and then go about differentiating.}

\item 
  \note{It's worthwhile to discuss efficiency and flexibility here.  For
    instance, students may want to decompose (b) as $f(g(t))$ with $f(t) =
    12e^t - e^2$ and $g(t) = t^3$; this works, but it might be simpler to
    first think about the fact that we have a difference, and that it's
    easy to take the derivative of the second piece of that difference. If
    you haven't done so already, this is a natural place to start having
    students work on problems at the board.}

  Differentiate each of the following functions. Try to be energy-efficient.
  Don't operate on automatic pilot.
  \begin{enumerate}

  \item
    \begin{problem}[\vfill]
      $h(x) = \ln (7x^2 + 3)$
    \end{problem}

    \begin{solution}
      We can decompose $h(x)$ as $f(g(x))$ with $g(x) = 7x^2 + 3$ and
      $f(x) = \ln x$.  Then, $g'(x) = 14 x$ and $f'(x) = \frac{1}{x}$, so
      $h'(x) = f'(g(x)) g'(x) = f'(7x^2 + 3) g'(x) = \boxed{\frac{1}{7x^2
          + 3} \cdot 14x}$ by the Chain Rule. 
    \end{solution}

  \item
    \begin{problem}[\vfill]
      $q(t) = e^{t^3} - 12 e^2 t$
    \end{problem}

    \begin{solution}
      We can decompose $e^{t^3}$ as $f(g(t))$ with $g(t) = t^3$ and $f(t)
      = e^t$.  Then, the Chain Rule says that $\frac{d}{dt} \left(e^{t^3}
      \right) = f'(g(t)) g'(t) = e^{t^3} \cdot 3t^2$, so $q'(t) =
      \boxed{3t^2 e^{t^3} - 12 e^2}$.

      Instead of writing out functions $f(t)$ and $g(t)$, we'll often
      write that we're going to use the Chain Rule simply by putting a red
      box around where we're going to use the Chain Rule ($e^{t^3}$ in
      this case) and then putting a blue box around the inner function
      $g(t)$: 
      \begin{align*}
        q'(t) &= \frac{d}{dt} \left(\fcolorbox{red}{white}{$\displaystyle
            e^ { \fcolorbox{blue}{white}{$\displaystyle t^ { 3} $}}
            $} - 12 e^2 t \right) \\ 
        &= e^ { t^ { 3} } \cdot 3 t^ { 2} -12 e^2
      \end{align*}
    \end{solution}

  \item
    \begin{problem}[\vfill\newpage]
      $h(x) = \pi^3 (7x + 3)^{10}$
    \end{problem}

    \begin{solution}
      Since $\pi^3$ is a constant,
      \begin{align*}
        h'(x) = {} & \pi ^3\cdot \frac{d}{dx} \left(\fcolorbox{red}{white}{$\displaystyle \fcolorbox{blue}{white}{$\displaystyle  \left(7 x+ 3\right)$}^ { 10} $}\right) \\
        = {} & \pi ^3\cdot 10 \left(7 x+ 3\right)^ { 9} \cdot 7
      \end{align*}
    \end{solution}

  \item
    \begin{problem}[\stretch{0.5}]
      $y = \displaystyle \frac{\sqrt{x^2 + 1}}{\pi}$.
    \end{problem}

    \begin{solution}
      We can rewrite $y$ as $\frac{1}{\pi} \sqrt{x^2 + 1}$; since
      $\frac{1}{\pi}$ is a constant,
      \begin{align*}
        y' = {} \frac{1}{\pi} \frac{d}{dx}
        \left(\fcolorbox{red}{white}{$\displaystyle
            \fcolorbox{blue}{white}{$\displaystyle  \left(x^2 + 1
              \right)$}^ { \frac{1}{2}} $}\right) = {} & \frac{1}{\pi} \cdot
        \frac{1}{2}
        \left(x^2 + 1 \right)^ { -\frac{1}{2}} \cdot
        \frac{d}{dx}\left(x^2 + 1 \right) \\ 
        = {} & \frac{1}{2\pi} \left(x^2 + 1 \right)^ { -\frac{1}{2}}
        \left(2x \right) 
      \end{align*}
    \end{solution}

  \item
    \begin{problem}[\vfill]
      $f(t) = \displaystyle \frac{5}{\sqrt[3]{1 + e^t}}$.  (Can you
      rewrite this so that you don't need to use the Quotient Rule?)
    \end{problem}

    \begin{solution}
      This can be rewritten as $f(t) = 5 (1 + e^t)^{-1/3}$.  Since $5$ is
      a constant, 
      \begin{align*}
        f'(t) = {} & 5\cdot \frac{d}{dt} \left(\fcolorbox{red}{white}{$\displaystyle \fcolorbox{blue}{white}{$\displaystyle  \left(1+ e^ { t} \right)$}^ { -\frac{1}{3}} $}\right) \\
        = {} & 5 \left(-\frac{1}{3} \left(1+ e^ { t} \right)^ { -\frac{4}{3}} \cdot \frac{d}{dt}\left(1+ e^ { t} \right)\right) \\
        = {} & -\frac{5}{3} \left(1+ e^ { t} \right)^ { -\frac{4}{3}}
        \cdot e^ { t} 
      \end{align*}
    \end{solution}

  \item
    \begin{problem}[\vfill]
      $h(x) = e^{x^2} \cdot \ln (7x^2 + 3)$
    \end{problem}

    \begin{solution}
      Here, the first rule we need is the Product Rule, because $h(x)$ is
      the product of $e^{x^2}$ and $\ln(7x^2 + 3)$; we'll show this in
      green by boxing $e^{x^2}$ and $\ln(7x^2 + 3)$ separately:
      \begin{align*}
        h'(x) = {} & \frac{d}{dx} \left(\fcolorbox{green}{white}{$\displaystyle  e^ { x^ { 2} } $}\cdot \fcolorbox{green}{white}{$\displaystyle  \ln\left(7 x^ { 2} + 3\right)$}\right) \\
        = {} & e^ { x^ { 2} } \cdot \frac{d}{dx} \left(\fcolorbox{red}{white}{$\displaystyle \ln\left(\fcolorbox{blue}{white}{$\displaystyle 7 x^ { 2} + 3$}\right)$}\right)+ \frac{d}{dx} \left(\fcolorbox{red}{white}{$\displaystyle e^ { \fcolorbox{blue}{white}{$\displaystyle x^ { 2} $}} $}\right)\cdot \ln\left(7 x^ { 2} + 3\right) \\
        = {} & e^ { x^ { 2} } \cdot \frac{1}{7 x^ { 2} + 3}\cdot 14 x+ e^
        { x^ { 2} } \cdot 2 x \cdot \ln\left(7 x^ { 2} + 3\right)
      \end{align*}
    \end{solution}

  \item
    \note{Preparation for logarithmic differentiation.}

    \begin{problem}[\vfill\newpage]
      $\displaystyle f(x) = \ln\left( \frac{x^2}{(x-2)(3x+5)}
      \right)$.  (This will be \textit{much} easier if you rewrite $f(x)$
      {before} differentiating.)
    \end{problem}

    \begin{solution}
      We can first use log rules to rewrite $f(x)$:
      \begin{align*}
        f(x) &= \ln x^2 - \ln [(x - 2) (3x + 5)] \\
        &= 2 \ln x - \left[ \ln(x - 2) + \ln (3x + 5)
        \right] \\
        &= 2 \ln x - \ln(x - 2) - \ln (3x + 5)
      \end{align*}
      Now let's differentiate, using the Chain Rule:
      \begin{align*}
        f'(x) = {} & 2 \cdot \frac{d}{dx} \left(\fcolorbox{red}{white}{$\displaystyle \ln\left(\fcolorbox{blue}{white}{$\displaystyle x$}\right)$}\right)-\frac{d}{dx} \left(\fcolorbox{red}{white}{$\displaystyle \ln\left(\fcolorbox{blue}{white}{$\displaystyle x-2$}\right)$}\right)-\frac{d}{dx} \left(\fcolorbox{red}{white}{$\displaystyle \ln\left(\fcolorbox{blue}{white}{$\displaystyle 3x+ 5$}\right)$}\right) \\
        = {} & 2\cdot \frac{1}{x}\cdot 1-\frac{1}{x-2}\cdot 1-\frac{1}{3x+ 5}\cdot 3
      \end{align*}
    \end{solution}

  \item
    \begin{problem}[\stretch{2}]
      $j(x) = [\ln (7x^2 + 3)]^5$
    \end{problem}

    \begin{solution}
      \begin{align*}
        j'(x) = {} & \frac{d}{dx}
        \left(\fcolorbox{red}{white}{$\displaystyle
            \fcolorbox{blue}{white}{$\displaystyle \ln\left(7 x^ { 2} +
                3\right)$}^ { 5} $}\right) \\
        = {} & 5 \left[ \ln\left(7 x^ { 2} + 3\right) \right]^ { 4} \cdot
        \frac{d}{dx} \left(\fcolorbox{red}{white}{$\displaystyle
            \ln\left(\fcolorbox{blue}{white}{$\displaystyle 7 x^ { 2} +
                3$}\right)$}\right) \\
        = {} & 5 \left[ \ln\left(7 x^ { 2} + 3\right) \right]^ { 4}
        \left(\frac{1}{7 x^ { 
              2} + 3}\cdot 14 x\right)
      \end{align*}
    \end{solution}

  \end{enumerate}

\item 

  \begin{enumerate}
  \item 
    \begin{problem}[\stretch{2}]
      How does the graph of $\ln(x + 5)$ relate to the graph of $\ln x$?
      Use the Chain Rule to differentiate $\ln (x + 5)$.  Is your answer
      consistent with the graph and what you know about the derivative of
      $\ln x$?
    \end{problem}

    \begin{solution}
      The graph of $\ln(x + 5)$ is simply the graph of $\ln x$ shifted 5
      units to the left.  Graphically, this means that the derivative of
      $\ln(x + 
      5)$ should be the derivative of $\ln x$ shifted left by 5.  That is,
      it should be $\frac{1}{x}$ shifted left by $5$, which has the
      formula $\frac{1}{x + 5}$. 

      Let's check.  By the Chain Rule,
      \[ \frac{d}{dx} \left(\fcolorbox{red}{white}{$\displaystyle
          \ln\left(\fcolorbox{blue}{white}{$\displaystyle x+
              5$}\right)$}\right) = \frac{1}{x+ 5}\cdot 1 = \frac{1}{x +
        5},\]
      as we expected.        
    \end{solution}

  \item
    \begin{problem}[\vfill]
      More generally, if $c$ is a constant, how does the graph of $f(x +
      c)$ relate to the graph of $f(x)$?  What does the Chain Rule tell
      you about the derivative of $f(x + c)$?
    \end{problem}

    \begin{solution}
      More generally, 
      the graph of $f(x + c)$ is simply the graph of $f(x)$ shifted $c$
      units to the left, so we expect the derivative of $f(x + c)$ (which
      we write as $\frac{d}{dx} f(x + c)$) to be $f'(x)$ shifted $c$ units
      to the left, or $f'(x + c)$. 

      By the Chain Rule,     
      \[ \frac{d}{dx} \left(\fcolorbox{red}{white}{$\displaystyle
          f\left(\fcolorbox{blue}{white}{$\displaystyle x+
              c$}\right)$}\right) = f'\left(x+ c\right)\cdot 1 = f'(x +
      c),\] as we expected. 
    \end{solution}    

  \end{enumerate}

\end{enumerate}


\end{document}